\subsection{セキュリティ管理と組織}
セキュリティ統治と管理は、個別の開発者の努力だけに依存すると弱い。組織の文化、行動、手順、責任分担に組み込まれているときに有効になる。プロジェクト管理者は、ソフトウェアセキュリティを単独の技術課題ではなく、組織全体の課題として扱う必要がある。

\subsubsection{能力成熟度モデル}
能力成熟度モデルとは、組織がどの程度一貫して、測定可能に、改善可能にセキュリティ活動を実行できるかを評価する考え方である。第13章では、システムセキュリティ工学能力成熟度モデル（SSE-CMM）が紹介される。

開発者、サービス提供者、システム統合者、システム管理者、セキュリティ専門家など、さまざまな役割がセキュリティ工学の方法を理解する必要がある。成熟度モデルは、リスク評価などを行う組織能力を測る道具として役立つ。

\subsubsection{情報セキュリティ管理システム}
情報セキュリティ管理システム（ISMS）は、組織の文脈で情報セキュリティを確立し、実施し、維持し、継続的に改善するための仕組みである。ISO/IEC 27001:2022 は、その要求事項を定める規格である。

ISMS は、技術に関わるセキュリティを管理する文書化された計画と考えられる。リスクを記録し、対応策を実施し、監視し、継続的に見直す。ISMS によって、新しいソフトウェアセキュリティ要求や、既存要求の変更が生じることもある。

\begin{notebox}{実務で見る点}セキュリティ要求は、顧客の希望だけから出るとは限らない。法令、規制、契約、監査、組織標準、ISMS のリスク評価からも生じる。\end{notebox}
\subsubsection{アジャイルにおけるソフトウェアセキュリティ実践}
アジャイルチームは、変化の速い開発の中でもセキュリティ実践を理解し、システムのセキュリティに責任を持つ必要がある。セキュリティ専門家も、変更を拒む立場ではなく、速く反復的に働き、リスクを小さな単位で扱う必要がある。

成功するアジャイルセキュリティには、セキュリティチームと開発者の関与、開発者が自分で実行できる支援、自動化、アジャイルチームの速度についていく俊敏性が必要である。セキュリティは「止める役割」ではなく、安全に進めるための「実現要因」であるべきである。
